
% !TeX root = ../main.tex
\begin{englishabstract}
    % \footnotetext{*The work was supported by the Foundation (foundation ID).}

    Indoor action recognition based on Wi-Fi Channel State Information (CSI) exploits channel variations caused by human motion to enable non-visual sensing, and has application value in smart homes, health monitoring, and environmental perception. Unlike offline recognition under a fixed label set, real-world systems often encounter progressively expanding action categories, incomplete retention of old-class data, and insufficient samples for newly introduced classes. Conventional training methods based on full static datasets are therefore difficult to adapt to continual update requirements. To address these issues, this thesis focuses on continuous category expansion in Wi-Fi CSI-based indoor action recognition and investigates two scenarios: class-incremental learning with relatively sufficient new-class samples and few-shot class-incremental learning with extremely limited new-class samples.

    For the class-incremental learning scenario, this thesis proposes a Replay-Expansion-Consolidation framework. Representative old-class samples are preserved through memory replay, new categories are incorporated by combining backbone freezing with classifier expansion, and old and new knowledge is consolidated through a teacher-student structure with bilateral knowledge distillation. This design alleviates catastrophic forgetting and the distribution imbalance between old and new classes. For the few-shot class-incremental learning scenario, this thesis proposes a learning method based on temporal enhancement and graph attention. Multi-scale temporal encoding and temporal attention fusion are used to enhance CSI temporal representations, an enhanced graph attention network is employed to model relations among class prototypes during incremental sessions, and a dual-classifier architecture is adopted to reduce old-new class bias and class-level overfitting.

    Experiments on the XRF and WiAR public Wi-Fi action recognition datasets show that the proposed class-incremental learning method maintains favorable stability during multi-stage category expansion. It achieves 73.57\% accuracy at the final incremental stage on XRF; under the WiAR-16 class-incremental protocol, it obtains an average accuracy of 77.33\% and a final-stage accuracy of 64.58\%. Under the 5-way 1-shot setting on XRF, the proposed few-shot class-incremental learning method achieves 74.24\% accuracy at the final session and an average accuracy of 78.48\%. On WiAR, it obtains an average accuracy of 75.30\% and a final-session accuracy of 61.81\%, outperforming CEC, FACT, C-FSCIL, DyCR, and other baselines overall. In summary, this thesis proposes continual learning methods for Wi-Fi CSI-based action recognition from the perspectives of class-incremental learning and few-shot class-incremental learning, providing a feasible solution for the continual updating of Wi-Fi sensing systems in dynamic scenarios.

    \englishkeywordstype{Wi-Fi CSI; Indoor action recognition; Continual learning; Class-incremental learning; Few-shot class-Incremental learning}{Application Research}

\end{englishabstract}
